\section{Parametri del sistema}\label{sec:par}

La tabella~\ref{tab:par-parametri} raccoglie i parametri del sistema descritto nei capitoli precedenti. Nell'ultima colonna «scelto» indica un valore fissato a tavolino, «tarato» un valore ricavato da dati simulati, «stimato» un valore preso da una fonte esterna o da un conto e non misurato. Nessun parametro è tarato su un sensore fisico; i valori marcati «non misurato» vanno fissati sulla camera usata. Le soglie di tempo valgono in secondi: il numero di fotogrammi è dato per il ciclo di $15\unita{Hz}$.

\begin{table}[H]\centering\footnotesize
\caption{Parametri del sistema (prima parte): sorgente, mappa, gruppi, filtro e ciclo di vita della traccia.}\label{tab:par-parametri}
\begin{tabularx}{\linewidth}{@{}p{1.25cm}Lp{3.5cm}Lp{1.2cm}@{}}
\toprule
Simbolo & Significato & Valore & Origine & Tipo\\
\midrule
$h$, $\theta$ & quota e inclinazione della camera & $1{,}2\unita{m}$; $15^\circ$ verso il basso & geometria d'osservazione (\S\ref{sec:geo}) & scelto\\
$\Delta$ & periodo del ciclo d'elaborazione & $1/15\unita{s}$ su un flusso a 30 fotogrammi al secondo & valore di progetto, non una prestazione misurata & scelto\\
--- & portata della mappa & $3\unita{m}$ & intervallo utile della camera; non misurato & scelto\\
--- & tolleranza sull'inclinazione & meno di $0{,}5^\circ$ & controesempio a $1^\circ$ (\S\ref{sec:lim-limiti}) & scelto\\
--- & autoverifica all'avvio & pixel validi sotto il $10\,\%$ entro 3 fotogrammi a sensore coperto & regola di progetto & scelto\\
--- & prove d'ingresso & oggetto di $10\unita{cm}$ a $2{,}5\unita{m}$ «occupato»; suolo vuoto per $30\unita{s}$ senza gruppi né tracce & regola di progetto & scelto\\
--- & cella della mappa polare & $1^\circ\times0{,}1\unita{m}$ & \S\ref{sec:map} & scelto\\
--- & quota d'ostacolo & da $0{,}07$ a $2{,}0\unita{m}$ sopra il suolo & \S\ref{sec:map} & scelto\\
--- & pixel minimi per «occupato» & 3 nel fotogramma corrente & \S\ref{sec:map} & scelto\\
--- & margine di lettura del suolo alla profondità attesa $z_0$ & $0{,}03+3\,k\,z_0^2$ (m); $k=0{,}004$ nominale, $0{,}0100$ al limite & modello del rumore del sensore; non misurato sul suolo & stimato\\
--- & condizione di «libero» & metà dei pixel dell'impronta della cella e di ognuna delle 8 vicine entro il margine & fissata dopo una prima simulazione & scelto\\
--- & memoria del «libero» & 0 fotogrammi & controesempi alle regole con memoria (\S\ref{sec:map}) & scelto\\
--- & inizio della frase sul suolo & $1{,}3\unita{m}$ & suolo visibile da $1{,}24\unita{m}$ sull'asse & scelto\\
--- & gruppi a celle forti (attivi solo dopo le prove d'ingresso) & celle con almeno $0{,}30$ del massimo della componente; $0{,}05\unita{m}$ e $0{,}004\unita{m^2}$ per le soglie attese & $0{,}30$ ricavato dopo i dati simulati; le altre fissate prima & tarato; scelto\\
$\sigma_z(d)$ & rumore di profondità & $1{,}5$ / $3{,}4$ / $6{,}0\unita{cm}$ a 2 / 3 / $4\unita{m}$ & dato del costruttore; non misurato & stimato\\
--- & soglia della regola «mobile» & $5\,\sigma\sqrt{2/m}$, $m=3$ fotogrammi per media, misure distanti al più $1\unita{s}$ & \S\ref{sec:sti} & scelto\\
$q$ & densità del rumore di processo & $0{,}2\unita{m^2/s^3}$ & prudente: $0{,}05$--$0{,}1$ su traiettorie pubbliche di pedoni & scelto\\
$\gamma$ & soglia di porta & quantile al $99\,\%$ & \S\ref{sec:sti} & scelto\\
--- & conferma della traccia & 3 misure su 5 & \S\ref{sec:sti} & scelto\\
--- & misure minime per la garanzia d'ingresso in ombra & 4 fotogrammi & con 3 la traccia manca nel $27{,}1\,\%$ dei casi, con 4 nel $5{,}2\,\%$ & tarato\\
--- & ingresso in ombra senza misura & dopo 2 cicli, con ultima misura o punto previsto entro $w$ dall'ombra & fissato dopo una prova sugli stessi semi & scelto\\
--- & durata di «non confermato» & $0{,}2\unita{s}$ & regola di progetto & scelto\\
--- & avviso senza traccia (attivo solo dopo le prove d'ingresso); distanza massima del riaggancio dall'ombra d'ingresso & $0{,}75\unita{m}$ & regola di progetto & scelto\\
--- & covarianza di misura al riaggancio & $R+(\text{larghezza}/2)^2$ & \S\ref{sec:sti} & scelto\\
\bottomrule
\end{tabularx}
\end{table}

\begin{table}[H]\centering\footnotesize
\caption{Parametri del sistema (seconda parte): occlusione, avviso, salute del sensore, uscita, profilo di chi porta il sensore.}\label{tab:par-parametri-due}
\begin{tabularx}{\linewidth}{@{}p{1.25cm}Lp{3.5cm}Lp{1.2cm}@{}}
\toprule
Simbolo & Significato & Valore & Origine & Tipo\\
\midrule
$w$ & fascia di guardia attorno all'ombra & $0{,}40\unita{m}$ & copre $1{,}13\unita{m/s}$ con $t_a=0{,}133\unita{s}$ (assunto), latenza $0{,}155\unita{s}$, tenuta $0{,}067\unita{s}$ & scelto\\
--- & fascia di ritaglio della regione sul suolo «libero» & $0{,}20\unita{m}$ & \S\ref{sec:occ} & scelto\\
$\alpha$ & livello della regione di previsione & $0{,}90$ & \S\ref{sec:occ} & scelto\\
--- & insieme di percorsi del predittore & 400, di cui 200 proseguono e 200 si fermano a un istante uniforme & \S\ref{sec:occ} & scelto\\
--- & peso di un percorso su suolo «libero» oltre la fascia di ritaglio & moltiplicato per $0{,}1$ a ogni ciclo & \S\ref{sec:occ} & scelto\\
--- & lisciature della densità & 3 & \S\ref{sec:occ} & tarato\\
--- & soglia e orizzonte di «previsto ampio» & regione oltre l'$80\,\%$ delle celle dell'ombra; $3\unita{s}$ & fissati dopo aver visto i dati simulati & scelto\\
--- & ingresso da fermo & velocità sotto $0{,}3\unita{m/s}$ all'ingresso o negli ultimi $2\unita{s}$ & \S\ref{sec:occ} & scelto\\
--- & punto previsto fuori dal non visto libero & fermo 8 cicli, poi al più $0{,}167\unita{m}$ per ciclo & \S\ref{sec:occ} & scelto\\
--- & libreria di percorsi ricavata da camminate di persone (non realizzata) & almeno 200 percorsi, due persone, tre andature & requisito di progetto; nelle prove la libreria è di 400 percorsi simulati (\S\ref{sec:occ}) & scelto\\
--- & rilascio di una cella dell'avviso & «libero» per 3 fotogrammi di fila e oltre $w$ dall'ombra corrente & \S\ref{sec:occ} & scelto\\
--- & fotogramma mancante & $0{,}13\unita{s}$ senza fotogramma nuovo, o tre identici & valore di progetto, non misurato; va fissato con $60\unita{s}$ di conteggio a vuoto e sotto carico & scelto\\
--- & pixel validi & sotto metà della taratura per $0{,}2\unita{s}$ & \S\ref{sec:lim-salute} & scelto\\
--- & unità inerziale; piano del suolo & $0{,}5^\circ$ o $5^\circ/\mathrm{s}$; $0{,}5^\circ$ & \S\ref{sec:lim-salute} & scelto\\
--- & segnale sonoro & $0{,}3\unita{s}$ per cambio di stato & \S\ref{sec:lim-uscita} & scelto\\
$\tau$ & latenza dall'evento all'uscita & $0{,}155\unita{s}$ (circa $0{,}131\unita{s}$ con un altro dato del costruttore) & somma di voci note (\S\ref{sec:lim-latenza}); non misurata & stimato\\
$t_{\mathrm{conf}}$ & tempo di conferma & $0{,}267\unita{s}$ & 4 periodi a $15\unita{Hz}$: valore usato nei conti per la conferma «3 misure su 5»; il legame fra i due numeri non è derivato in questo documento & scelto\\
$C$ & margine della distanza d'avviso & $0{,}95\unita{m}$ & $0{,}85\unita{m}$ attribuiti a una norma, non riscontrati, più $0{,}10\unita{m}$ & stimato\\
$T_r$, $v_e$, $a_f$ & tempo di reazione, velocità propria, frenata del portatore & per profilo (tabella~\ref{tab:lim-portatore}) & ipotesi di progetto o valori riportati & scelto\\
--- & soglia di fiducia del classificatore; coerenza cinematica & $0{,}7$; 4 cicli su 6 & ipotesi; il classificatore non è realizzato (\S\ref{sec:sti}) & scelto\\
\bottomrule
\end{tabularx}
\end{table}

\paragraph{Parametri a lunga portata.} La tabella~\ref{tab:par-lunga} raccoglie i parametri del \S\ref{sec:lp}; i valori dei sensori stanno nella tabella~\ref{tab:lp-sensori}. Nessuno è tarato su un sensore fisico.

\begin{table}[H]\centering\footnotesize
\caption{Parametri del sistema a lunga portata: «libero» per classe, inseguimento, avviso, eventi.}\label{tab:par-lunga}
\begin{tabularx}{\linewidth}{@{}p{1.25cm}Lp{3.5cm}Lp{1.2cm}@{}}
\toprule
Simbolo & Significato & Valore & Origine & Tipo\\
\midrule
--- & lato della cella polare & $100\unita{m}$ & radar di riferimento delle prove del «libero» & scelto\\
$P_c$ & soglia del «libero» per classe & $0{,}99$ & ipotesi & scelto\\
$W_c$ & finestra di cumulo delle passate & $\min(t_{\max},\ell/v_{\max,c})$, al più 5 passate & \S\ref{sec:lp-libero} & scelto\\
$\rho_c$ & correlazione fra passate & 1 ($n_{\mathrm{eff}}=1$) & prudente (Proposizione~\ref{prop:lp-correlate}) & scelto\\
$P_d$ & probabilità di rilevamento & Swerling~1; $P_d=0{,}5$ alla portata di scheda, $P_{fa}=10^{-6}$ & ipotesi di modello & stimato\\
$\sigma$ & sezione radar del multirotore piccolo; del drone sotto $250\unita{g}$ & $0{,}01\unita{m^2}$; $1{,}6\cdot10^{-4}\unita{m^2}$ & ipotesi; il secondo da un conto sulla scheda & stimato\\
--- & sfondo misurato per accertare la linea di vista & entro $2\unita{m}$ dall'atteso & \S\ref{sec:lp-libero} & scelto\\
--- & porta di associazione; porta del dubbio & $d^2\le16{,}27$; $16{,}27<d^2\le30{,}66$, terza misura dubbia in $5\unita{s}$ apre & chi quadro con 3 gradi di libertà & scelto\\
--- & conferma della traccia & 2 misure su 3; tentativa tolta dopo $3\unita{s}$ & \S\ref{sec:lp-inseguimento} & scelto\\
--- & misure fuori sequenza & scartate se più vecchie di $0{,}5\unita{s}$ & \S\ref{sec:lp-inseguimento} & scelto\\
$v_{\max}$, $a_{\max}$ & inviluppo del sito sorvegliato & multirotore veloce $45\unita{m/s}$, $25\unita{m/s^2}$ & ipotesi senza fonte & scelto\\
$v_{\max}$ & inviluppo degli aerei in volo & $210\unita{m/s}$ & decolli reali fino a $205\unita{m/s}$ (\S\ref{sec:lp-traiettorie}) & stimato\\
$w$ & fascia dell'insieme di possibile presenza & $50\unita{m}$ & ipotesi & scelto\\
$\alpha$ & livello della zona prevista & $0{,}90$, orizzonti 1, 2, 5, 10, $30\unita{s}$ & quantile conforme su traiettorie & tarato\\
--- & «previsto ampio» & zona oltre $0{,}8$ del raggio dell'avviso & \S\ref{sec:lp-presenza} & scelto\\
$\tau_{\mathrm{lat}}$, $T_r$, $T_{\mathrm{man}}$ & latenza, reazione e messa in sicurezza dell'operatore di sito & $1$, $5$, $30\unita{s}$; margine $50\unita{m}$ & ipotesi & scelto\\
--- & innesco di «separazione» & almeno 3 ritorni nuovi confermati entro $2\unita{s}$ & \S\ref{sec:lp-eventi} & scelto\\
$\Delta v$, $W$, $W_z$ & velocità dei frammenti; vento orizzontale e verticale (collisione, valori nominali) & $40$, $10$, $1\unita{m/s}$ & ipotesi & scelto\\
--- & frammento minimo & $5\unita{cm}$, $2\cdot10^{-4}\unita{m^2}$ & ipotesi & scelto\\
--- & nube & 3 cicli, più $10\unita{m}$; tolta dopo $5\unita{s}$ di sfondo visibile & \S\ref{sec:lp-eventi} & scelto\\
--- & abbagliamento & pixel saturi $>5\,\%$ (visibile), $>2\,\%$ (termica); ritorni LiDAR $-30\,\%$; fondo radar $+6\unita{dB}$; entro 3 cicli & ipotesi & scelto\\
\bottomrule
\end{tabularx}
\end{table}
